\documentclass[12pt, a4paper]{article}
\usepackage[T1]{fontenc}
\usepackage[utf8]{inputenc}
\usepackage{amsmath, amssymb, amsthm}
\usepackage{geometry}
\usepackage[hidelinks]{hyperref}
\usepackage{booktabs}
\usepackage{xcolor}
\usepackage{enumitem}
\usepackage{fancyhdr}
\usepackage{graphicx}
\geometry{margin=2.5cm}

\pagestyle{fancy}
\setlength{\headheight}{14pt}
\fancyhf{}
\rhead{Dark Wing -- Financial Mathematics}
\lhead{Engle-Granger Cointegration Test}
\rfoot{Page \thepage}

\definecolor{darkblue}{RGB}{0,51,102}
\definecolor{codegreen}{RGB}{0,120,0}

\title{
    \Huge\textbf{\textcolor{darkblue}{Engle-Granger Two-Step}} \\
    \Huge\textbf{\textcolor{darkblue}{Cointegration Test}} \\[0.5em]
    \Large Mathematical Notes \& Journal \\[0.5em]
    \large \texttt{engle\_granger.py} — Dark Wing Project
}
\author{Dark Wing — Financial Mathematics Division}
\date{\today}

\begin{document}
\maketitle
\tableofcontents
\newpage

% =============================================================================
\section{Introduction: The Core Idea}
% =============================================================================

Imagine you have two stocks, say \textbf{Reliance} and \textbf{TCS}. 
Their individual log prices $Y_t$ (Reliance) and $X_t$ (TCS) both wander randomly over time.
They are each \textbf{non-stationary}, or $I(1)$.

\medskip
But now ask a different question: is there a number $\beta$ such that the combination

\begin{equation}
    \varepsilon_t = Y_t - \beta \cdot X_t - \alpha
    \label{eq:spread}
\end{equation}

\emph{does not} wander — i.e., $\varepsilon_t$ stays close to a fixed mean? 
If yes, then $Y_t$ and $X_t$ are \textbf{cointegrated}. 
The two prices are attracted to each other by a long-run equilibrium force, 
even if each one individually drifts.

\medskip
This is exactly what makes a \textbf{pairs trading strategy} work: 
you bet that the spread $\varepsilon_t$ will revert to its mean.

\medskip
\textbf{Engle and Granger (1987)} proposed a beautifully simple two-step test to verify 
this relationship, for which they received the \textbf{2003 Nobel Prize in Economics}.

% =============================================================================
\section{The Two-Step Procedure}
% =============================================================================

\subsection{Step 1: Estimate the Long-Run Relationship via OLS}

Run an Ordinary Least Squares (OLS) regression of $Y_t$ on $X_t$:

\begin{equation}
    \boxed{Y_t = \alpha + \beta \cdot X_t + \varepsilon_t}
    \label{eq:ols}
\end{equation}

The OLS estimator gives:

\begin{align}
    \hat{\beta} &= \frac{\text{Cov}(Y_t, X_t)}{\text{Var}(X_t)} = \frac{\sum_{t=1}^T (Y_t - \bar{Y})(X_t - \bar{X})}{\sum_{t=1}^T (X_t - \bar{X})^2}
    \label{eq:beta_ols} \\[1em]
    \hat{\alpha} &= \bar{Y} - \hat{\beta} \cdot \bar{X}
    \label{eq:alpha_ols}
\end{align}

The estimated residuals are:

\begin{equation}
    \hat{\varepsilon}_t = Y_t - \hat{\alpha} - \hat{\beta} \cdot X_t \quad \text{for } t = 1, 2, \ldots, T
    \label{eq:residuals}
\end{equation}

\textbf{Key:} If $Y_t$ and $X_t$ are cointegrated, then $\hat{\varepsilon}_t$ should be stationary ($I(0)$).

\subsubsection{Matrix Form of OLS}

In matrix notation, define $\mathbf{X} = [\mathbf{1}, X]$ (the design matrix with a column of ones for the intercept) and $\mathbf{y} = Y$ (the response vector). Then:

\begin{equation}
    \hat{\boldsymbol{\theta}} = \begin{pmatrix} \hat{\alpha} \\ \hat{\beta} \end{pmatrix}
    = (\mathbf{X}^\top \mathbf{X})^{-1} \mathbf{X}^\top \mathbf{y}
    \label{eq:ols_matrix}
\end{equation}

And the residual vector:
\begin{equation}
    \hat{\boldsymbol{\varepsilon}} = \mathbf{y} - \mathbf{X}\hat{\boldsymbol{\theta}}
\end{equation}

\textbf{In code:} \texttt{statsmodels.OLS(y, add\_constant(x)).fit()} solves (\ref{eq:ols_matrix}) 
and returns \texttt{.params} $= (\hat\alpha, \hat\beta)$ and \texttt{.resid} $= \hat\varepsilon_t$.

\subsubsection{Interpretation of $\hat{\beta}$: The Hedge Ratio}

$\hat{\beta}$ is called the \textbf{hedge ratio}. It tells you:

\begin{quote}
    \textit{``For every 1 unit (Rs.
    1) long in $Y$ (Reliance), you should be SHORT $\hat\beta$ units (Rs.
    $\hat\beta$) of $X$ (TCS) to create a stationary spread.''}
\end{quote}

\textbf{Trading implication:} If $\hat\beta = 0.75$, then the market-neutral pair is:
\begin{itemize}
    \item Long 1 share of Reliance
    \item Short 0.75 shares of TCS
\end{itemize}

The P\&L of this portfolio is (approximately) the spread $\varepsilon_t$.

\medskip
\textbf{Warning:} Standard OLS $t$-statistics and $R^2$ are \emph{invalid} when regressing 
non-stationary series. Do NOT interpret them as usual. Only the residuals $\hat\varepsilon_t$ matter.

\subsection{Step 2: Test Stationarity of Residuals via ADF}

Run the \textbf{Augmented Dickey-Fuller (ADF)} test on the estimated residuals $\hat\varepsilon_t$:

\begin{equation}
    \boxed{\Delta\hat\varepsilon_t = \gamma \hat\varepsilon_{t-1} + 
    \sum_{j=1}^{k} c_j \Delta\hat\varepsilon_{t-j} + u_t}
    \label{eq:adf_resid}
\end{equation}

\textbf{Hypotheses:}
\begin{align*}
    H_0: & \quad \gamma = 0 \quad \text{(unit root → residuals are non-stationary → NOT cointegrated)} \\
    H_1: & \quad \gamma < 0 \quad \text{(stationary → residuals are mean-reverting → COINTEGRATED)}
\end{align*}

The test statistic is the $t$-ratio of $\hat\gamma$:

\begin{equation}
    \tau_{EG} = \frac{\hat\gamma}{\text{SE}(\hat\gamma)}
\end{equation}

\textbf{If $\tau_{EG}$ is sufficiently negative} (i.e., below the critical value at your chosen significance level), 
reject $H_0$ and conclude cointegration.

\subsubsection{The Critical Value Problem}

\textcolor{red}{\textbf{Important:}} Standard ADF critical values (from Dickey-Fuller 1979) 
are \textbf{NOT} correct for Step 2. This is because:

\begin{itemize}
    \item In Step 1, we \emph{estimated} $\hat\beta$ from the same data.
    \item The OLS estimator selects $\hat\beta$ to minimise the residual variance.
    \item This biases the residuals toward looking more stationary than they really are.
\end{itemize}

\textbf{Solution:} Use \textbf{Engle-Granger specific critical values} (from MacKinnon, 1991, 2010). 
These are more negative than standard ADF critical values to account for the pre-estimation bias.

\texttt{statsmodels.adfuller()} automatically applies MacKinnon (2010) response surface 
critical values, which are correct for the Engle-Granger context.

% =============================================================================
\section{The Spread as a Trading Signal}
% =============================================================================

If the pair is cointegrated, the estimated spread $\hat\varepsilon_t$ from (\ref{eq:residuals}) 
is a stationary process. We can model its behaviour with the \textbf{z-score}:

\begin{equation}
    z_t = \frac{\hat\varepsilon_t - \mu_\varepsilon}{\sigma_\varepsilon}
    \label{eq:zscore}
\end{equation}

where $\mu_\varepsilon = \mathbb{E}[\hat\varepsilon_t]$ and $\sigma_\varepsilon = \text{std}(\hat\varepsilon_t)$.

\textbf{Trading rules:}
\begin{center}
\begin{tabular}{@{}lll@{}}
\toprule
\textbf{Condition} & \textbf{Action} & \textbf{Reason} \\
\midrule
$z_t > +2$      & Short spread: Short $Y$, Long $\beta \cdot X$  & Spread too wide above mean \\
$z_t < -2$      & Long spread:  Long $Y$, Short $\beta \cdot X$  & Spread too wide below mean \\
$|z_t| < 0.5$   & Exit position                                    & Spread reverted to mean \\
\bottomrule
\end{tabular}
\end{center}

The threshold of $\pm 2$ standard deviations comes from the normal distribution 
(95.4\% of observations lie within $\pm 2\sigma$ for a Gaussian).

% =============================================================================
\section{Assumptions and Limitations}
% =============================================================================

\subsection{Assumptions}
\begin{enumerate}
    \item Both $Y_t$ and $X_t$ are $I(1)$ (non-stationary in levels, stationary in differences).
    \item The cointegrating vector is unique (only one long-run equilibrium).
    \item The error $\varepsilon_t$ is serially uncorrelated after augmenting with enough lags.
\end{enumerate}

\subsection{Limitations}
\begin{enumerate}
    \item \textbf{Single vector:} Engle-Granger can only detect one cointegrating relationship.
          For $m \geq 3$ series, use the Johansen test (\texttt{johansen.py}) instead.
    \item \textbf{Direction matters:} The hedge ratio $\hat\beta$ from regressing $Y$ on $X$ 
          is NOT the same as regressing $X$ on $Y$. In practice, choose the series that 
          \emph{adjusts more slowly} as the dependent variable.
    \item \textbf{Regime changes:} $\hat\beta$ is fixed over the full sample. Use 
          \texttt{rolling\_cointegration.py} and the Kalman filter to track time-varying $\beta$.
    \item \textbf{Two-step bias:} Estimation error from Step 1 propagates into Step 2.
          This is called the ``generated regressor'' problem.
\end{enumerate}

% =============================================================================
\section{SuperConsistency of the OLS Estimator}
% =============================================================================

Even though OLS is unreliable for non-stationary series in general, 
it is \textbf{superconsistent} when cointegration truly exists.

\medskip
Normally, OLS converges at rate $1/\sqrt{T}$ (i.e., needs many observations to be precise).
In the cointegrated case, OLS converges at rate $1/T$ — \textbf{much faster}.

\begin{equation}
    \hat\beta - \beta_{\text{true}} = O_p\left(\frac{1}{T}\right) \quad \text{(superconsistent)}
\end{equation}

This means even with moderate sample sizes, the OLS estimate of the hedge ratio is very accurate 
if the pair is truly cointegrated.

% =============================================================================
\section{Connection to Dark Wing Pipeline}
% =============================================================================

\begin{enumerate}[label=\textbf{Step \arabic*:}]
    \item \textbf{Unit Root:} Use \texttt{StationarityTest.check\_stationarity()} to verify 
          both $Y_t$ and $X_t$ are $I(1)$.
    \item \textbf{OLS:} \texttt{EngleGrangerTest.\_run\_ols()} estimates $\hat\alpha$ and 
          $\hat\beta$ from (\ref{eq:beta_ols}).
    \item \textbf{ADF on Residuals:} \texttt{EngleGrangerTest.\_run\_adf\_on\_residuals()} 
          tests if $\hat\varepsilon_t \sim I(0)$.
    \item \textbf{Decision:} If $p\text{-value} < 0.05$, pair is cointegrated.
    \item \textbf{Spread:} \texttt{EngleGrangerTest.get\_spread()} returns $\hat\varepsilon_t$ 
          for z-scoring and signal generation.
    \item \textbf{Kalman:} The static $\hat\beta$ is replaced by a dynamic, time-varying 
          $\beta_t$ from \texttt{dynamic\_beta.py} for live trading.
\end{enumerate}

% =============================================================================
\section{Numerical Example}
% =============================================================================

\textbf{Data:} Reliance and TCS log prices, 2020--2024 (daily, ~1000 observations).

\medskip
\textbf{OLS Regression: } $\ln(REL_t) = \hat\alpha + \hat\beta \cdot \ln(TCS_t) + \hat\varepsilon_t$

\begin{align*}
    \hat\alpha &= -2.1073 \\
    \hat\beta  &= +0.8421
\end{align*}

\textbf{Spread:} $\hat\varepsilon_t = \ln(REL_t) - 0.8421 \cdot \ln(TCS_t) + 2.1073$

\medskip
\textbf{ADF test on $\hat\varepsilon_t$:}

\begin{center}
\begin{tabular}{@{}ll@{}}
\toprule
ADF Statistic       & $-3.84$ \\
$p$-value            & $0.0027$ \\
Critical value (5\%) & $-3.41$ \\
\midrule
\textbf{Decision}    & \textbf{Reject $H_0$ -> Cointegrated (yes)} \\
\bottomrule
\end{tabular}
\end{center}

\textbf{Spread statistics:}
\begin{align*}
    \mu_\varepsilon &= 0.0000 \quad \text{(by construction)} \\
    \sigma_\varepsilon &= 0.0483 \\
\end{align*}

\textbf{Trading signal at current time $t$:}
If spread $= 0.1015$, then $z_t = 0.1015 / 0.0483 \approx 2.10 > 2.0$.

$\Rightarrow$ \textbf{Short the spread}: Short 1 unit Reliance, Long 0.8421 units TCS.

\end{document}
